\chapter{Generating Functions}

The previous chapter showed that the distribution of a sum of independent random
variables is given by convolution. Calculating repeated convolutions by hand,
however, quickly becomes unmanageable.

Generating functions solve this by encoding an entire probability distribution
into a single function, transforming convolution into ordinary multiplication
and moments into derivatives. We study three versions: the \emph{probability
generating function} (PGF) for discrete counts, the \emph{moment generating
function} (MGF) for general random variables, and the \emph{characteristic
function} (CF), which is guaranteed to exist for every distribution.

\chapterline

\section{The Idea: Hang Probabilities on a Function}
\label{ch18:idea}

Let us start with a single die and a deliberately small question, so the trick
is visible before it is useful.

A fair die produces a value in $\set{1,2,3,4,5,6}$, each with probability
$1/6$. Here is an unusual way to write that distribution down. Introduce a
bookkeeping symbol $s$ and form the polynomial
\[
  G(s) \;=\; \tfrac16 s^1 + \tfrac16 s^2 + \tfrac16 s^3
            + \tfrac16 s^4 + \tfrac16 s^5 + \tfrac16 s^6 .
\]
Read it like a coat rack. The coefficient hanging on the hook labelled $s^k$ is
the probability that the die shows $k$. Nothing has been added or lost --- the
six numbers $1/6$ are all still there --- but now they ride on a single
function of $s$. The exponent records the \emph{value}, the coefficient records
the \emph{probability}.

\begin{intuitionbox}
A distribution on $\set{0,1,2,\dots}$ is just a sequence of numbers
$p_0,p_1,p_2,\dots$. A generating function is the trick of bundling a whole
sequence into one object by using powers of a dummy variable as labelled
shelves: $G(s)=p_0 + p_1 s + p_2 s^2 + \cdots$. The shelf $s^k$ holds $p_k$.
What makes this worth doing is that natural operations on \emph{sequences}
(like convolution) become natural operations on \emph{functions} (like
multiplication).
\end{intuitionbox}

Why would you ever want this? Watch what happens when you roll \emph{two}
independent dice and ask for the distribution of the total. The total is $k$
when the first die shows $i$ and the second shows $k-i$, for every way that can
happen, so the probability of total $k$ is the convolution
\[
  \P(\text{sum}=k) \;=\; \sum_{i} \P(\text{first}=i)\,\P(\text{second}=k-i).
\]
Now look at the product of the two coat racks. Multiplying $G(s)\cdot G(s)$ and
collecting the coefficient of $s^k$, you pick up exactly
$\sum_i (\text{coeff of }s^i)(\text{coeff of }s^{k-i})$ --- because $s^i \cdot
s^{k-i}=s^k$. \emph{That is the convolution sum, performed for free by the
exponent arithmetic.} The exponents add precisely because random variables add.
Convolution of probabilities and multiplication of polynomials are the same
operation seen from two sides.

\begin{happeningbox}
The whole subject lives in one observation: $s^i\cdot s^j = s^{i+j}$. When you
multiply two generating functions, every cross term $s^i$-from-the-first times
$s^j$-from-the-second lands on the shelf $s^{i+j}$, carrying the product of its
two probabilities. Summing all the ways to land on shelf $k$ \emph{is} the
convolution that gives $\P(X+Y=k)$. So ``add independent random variables''
$\longleftrightarrow$ ``multiply their generating functions.'' Everything else
in this chapter is a consequence and a refinement of that single line.
Figure~\ref{ch18:fig-coatrack-conv} shows this multiplication in action for two dice.
\end{happeningbox}

\begin{figure}[H]
\centering
\begin{tikzpicture}[scale=0.92]
  \begin{scope}[shift={(0,0)}]
    \node[anchor=south, font=\small\bfseries] at (2.4, 4.5) {1. Outcome Pairs $(i,j)$};
    \foreach \i in {1,...,6} {
      \foreach \j in {1,...,6} {
        \pgfmathtruncatemacro{\k}{\i+\j}
        \ifnum\k=7
          \fill[quotientorange!35, draw=quotientorange!70, rounded corners=1pt] 
            ({\i*0.65-0.28}, {\j*0.65-0.28}) rectangle ({\i*0.65+0.28}, {\j*0.65+0.28});
          \node[font=\tiny\bfseries, quotientorange!90!black] at ({\i*0.65}, {\j*0.65}) {$s^7$};
        \else
          \fill[gray!6, draw=gray!25, rounded corners=1pt] 
            ({\i*0.65-0.28}, {\j*0.65-0.28}) rectangle ({\i*0.65+0.28}, {\j*0.65+0.28});
          \node[font=\tiny, gray!60] at ({\i*0.65}, {\j*0.65}) {$s^{\k}$};
        \fi
      }
    }
    \draw[->, thick, gray!70] (0.2, 0.35) -- (4.4, 0.35) node[right, font=\scriptsize] {Die 1 ($i$)};
    \draw[->, thick, gray!70] (0.35, 0.2) -- (0.35, 4.2) node[above, font=\scriptsize] {Die 2 ($j$)};
    \foreach \i in {1,...,6} {
      \node[font=\tiny, gray!80] at ({\i*0.65}, 0.15) {$\i$};
      \node[font=\tiny, gray!80] at (0.15, {\i*0.65}) {$\i$};
    }
    \draw[quotientorange, thick, dashed] (0.65*1, 0.65*6) -- (0.65*6, 0.65*1)
      node[above right=2pt, font=\scriptsize\bfseries, quotientorange!90!black] {$i+j = 7$};
  \end{scope}

  \begin{scope}[shift={(4.7, 2.1)}]
    \draw[->, >=stealth, very thick, sourceblue] (0, 0) -- (1.4, 0);
    \node[align=center, font=\scriptsize\bfseries, sourceblue, anchor=south] at (0.7, 0.1) {Multiply\\$s^i \cdot s^j = s^{i+j}$};
    \node[font=\tiny, gray!70, anchor=north] at (0.7, -0.1) {Collect like terms};
  \end{scope}

  \begin{scope}[shift={(6.9, 0)}]
    \node[anchor=south, font=\small\bfseries] at (3.5, 4.5) {2. Coefficients of $G(s)^2$ (Sum PMF)};
    \draw[->, thick, gray!70] (0, 0.35) -- (6.6, 0.35) node[right, font=\scriptsize] {power of $s$};
    \draw[->, thick, gray!70] (0.3, 0.2) -- (0.3, 4.2) node[above left, font=\scriptsize] {$\P(X+Y=k)$};

    \foreach \k/\h in {2/1, 3/2, 4/3, 5/4, 6/5, 8/5, 9/4, 10/3, 11/2, 12/1} {
      \pgfmathsetmacro{\xpos}{0.5 * (\k - 2) + 0.8}
      \pgfmathsetmacro{\yht}{0.48 * \h}
      \fill[sourceblue!45, draw=sourceblue] ({\xpos-0.18}, 0.35) rectangle ({\xpos+0.18}, {0.35+\yht});
      \node[font=\tiny, gray!80] at (\xpos, 0.15) {$s^{\k}$};
    }
    \pgfmathsetmacro{\xseven}{0.5 * (7 - 2) + 0.8}
    \fill[quotientorange!55, draw=quotientorange!90, very thick] 
      ({\xseven-0.18}, 0.35) rectangle ({\xseven+0.18}, {0.35 + 0.48*6});
    \node[font=\tiny\bfseries, quotientorange!90!black] at (\xseven, 0.15) {$s^7$};
    \node[font=\scriptsize\bfseries, quotientorange!90!black, anchor=south] 
      at (\xseven, {0.35 + 0.48*6}) {$\frac{6}{36}$};

    \node[anchor=west, align=center, font=\scriptsize, draw=quotientorange!70, fill=white, rounded corners=2pt, inner sep=2pt]
      at (3.8, 3.6) {6 terms on diagonal $i+j=7$\\[1pt]combine into $\frac{6}{36}s^7$};
  \end{scope}
\end{tikzpicture}
\caption{How polynomial multiplication automates convolution. (Left) Rolling two dice produces $36$ equally likely outcome pairs $(i,j)$, each carrying the term $\frac{1}{36}s^i s^j = \frac{1}{36}s^{i+j}$. Exponents add ($s^i \cdot s^j = s^{i+j}$), so all pairs lying on the diagonal slice $i+j=k$ share the exact same power $s^k$. (Right) Multiplying the generating functions $G(s)\cdot G(s)$ automatically gathers all terms on diagonal $k$ into the coefficient of $s^k$, producing the triangular sum distribution without any manual summation.}
\label{ch18:fig-coatrack-conv}
\end{figure}

\section{The Probability Generating Function}
\label{ch18:pgf}

Let us make the coat rack official, for random variables that take values in
$\set{0,1,2,\dots}$.

\begin{definition}[Probability generating function]
\label{ch18:def-pgf}
Let $X$ be a random variable taking values in the nonnegative integers, with
pmf $p_k=\P(X=k)$. Its \emph{probability generating function} is
\[
  G_X(s) \;=\; \E\!\left[s^X\right] \;=\; \sum_{k=0}^{\infty} p_k\, s^k ,
  \qquad \abs{s}\le 1 .
\]
\end{definition}

The two faces of the definition --- ``$\E[s^X]$'' and ``$\sum p_k s^k$'' --- are
worth holding side by side. The first is compact and is how you will usually
\emph{compute} it; the second is the coat rack and is how you will usually
\emph{read} it. They agree because $\E[s^X]=\sum_k s^k\,\P(X=k)$ is just the
definition of expectation of the function $s\mapsto s^X$ (this is LOTUS, the law
of the unconscious statistician, from the expectation chapter).

The series converges at least for $\abs{s}\le 1$, because there the terms are
dominated by the probabilities: $\sum_k p_k \abs{s}^k \le \sum_k p_k = 1$.
In particular $G_X(1)=\sum_k p_k = 1$ always, which is a handy sanity check ---
a correct PGF must pass through the point $(1,1)$.

\begin{recoverybox}
\textbf{Mental check: reading the edges of a PGF.}
Every valid PGF comes with two instant sanity checks:
\begin{itemize}
  \item At $s=1$: $G_X(1) = \sum_{k=0}^\infty p_k = 1$. The probabilities must sum to $1$.
  \item At $s=0$: $G_X(0) = p_0 \cdot 0^0 + p_1 \cdot 0^1 + p_2 \cdot 0^2 + \cdots = p_0 = \P(X=0)$. The value at the origin is always the probability of zero!
\end{itemize}
For example, for the Poisson PGF $G(s) = e^{\lambda(s-1)}$, checking $s=0$ gives $G(0) = e^{-\lambda}$, exactly $\P(X=0)$.
\end{recoverybox}

\begin{example}[The usual suspects]
\label{ch18:ex-pgfs}
A few PGFs you should be able to write on sight.
\begin{itemize}
  \item \textbf{Bernoulli}$(p)$: $X$ is $1$ w.p.\ $p$ and $0$ w.p.\ $q=1-p$, so
        \[
          G_X(s) = q\,s^0 + p\,s^1 = q + p s .
        \]
  \item \textbf{Binomial}$(n,p)$: think of $X=\sum_{i=1}^n X_i$ as a sum of $n$
        independent Bernoulli$(p)$'s. We will see in a moment that PGFs of
        independent sums multiply, so
        \[
          G_X(s) = (q+ps)^n .
        \]
        (Notice the binomial theorem expands this into
        $\sum_k \binom{n}{k}p^k q^{n-k}s^k$ --- the pmf, recovered as
        coefficients.)
  \item \textbf{Geometric}$(p)$ (number of failures before the first success,
        $k=0,1,2,\dots$): $p_k = q^k p$, so
        \[
          G_X(s) = \sum_{k=0}^\infty q^k p\, s^k
                 = p\sum_{k=0}^\infty (qs)^k = \frac{p}{1-qs},
          \qquad \abs{s}<1/q .
        \]
  \item \textbf{Poisson}$(\lambda)$: $p_k = e^{-\lambda}\lambda^k/k!$, so
        \[
          G_X(s) = \sum_{k=0}^\infty e^{-\lambda}\frac{\lambda^k}{k!}s^k
                 = e^{-\lambda}\sum_{k=0}^\infty \frac{(\lambda s)^k}{k!}
                 = e^{-\lambda}e^{\lambda s} = e^{\lambda(s-1)} .
        \]
        The Taylor series of $e^x$ does the summing for you, as it so often
        does with the Poisson.
\end{itemize}
\end{example}

\subsection{The multiplication theorem}

Now the property the whole construction was built to have.

\begin{theorem}[PGF of an independent sum]
\label{ch18:thm-pgf-product}
If $X$ and $Y$ are independent nonnegative-integer random variables, then
\[
  G_{X+Y}(s) \;=\; G_X(s)\,G_Y(s).
\]
More generally, if $X_1,\dots,X_n$ are independent, then $G_{X_1+\cdots+X_n}(s)
= \prod_{i=1}^n G_{X_i}(s)$.
\end{theorem}

\begin{whybox}
The proof is one line, and it is the most important line in the chapter:
\[
  G_{X+Y}(s) = \E\!\left[s^{X+Y}\right]
             = \E\!\left[s^{X}s^{Y}\right]
             \stackrel{\indep}{=} \E\!\left[s^{X}\right]\E\!\left[s^{Y}\right]
             = G_X(s)\,G_Y(s).
\]
The single starred step uses independence: the expectation of a product of
independent quantities is the product of their expectations. Everything else is
exponent arithmetic, $s^{X+Y}=s^X s^Y$. This is the convolution-becomes-product
fact from \S\ref{ch18:idea}, now stated cleanly in the language of expectation.
The $n$-fold version follows by induction.
\end{whybox}

\begin{example}[Sum of independent Poissons via PGFs]
\label{ch18:ex-pois-sum}
Let $X\sim\Pois(\lambda)$ and $Y\sim\Pois(\mu)$ be independent. By
Theorem~\ref{ch18:thm-pgf-product},
\[
  G_{X+Y}(s) = e^{\lambda(s-1)}\,e^{\mu(s-1)} = e^{(\lambda+\mu)(s-1)} .
\]
But that last expression is exactly the PGF of a $\Pois(\lambda+\mu)$ variable.
Once we know (next section) that the PGF determines the distribution, we may
conclude $X+Y\sim\Pois(\lambda+\mu)$ --- no convolution sum, no
$\sum_{j}\binom{k}{j}\lambda^j\mu^{k-j}/k!$ to wrestle with. The exponents just
added. This is the payoff promised in the chapter opening, delivered in two
lines.
\end{example}

\begin{figure}[H]
\centering
\begin{tikzpicture}
\begin{axis}[
  width=13cm, height=5.6cm,
  ybar, bar width=4pt,
  ymin=0, ymax=0.25,
  xmin=-0.5, xmax=20.5,
  xlabel={value $k$}, ylabel={probability},
  title={Adding independent Poissons just adds the rate},
  title style={font=\small},
  xtick={0,4,8,12,16,20},
  ytick={0,0.05,0.10,0.15},
  scaled y ticks=false,
  yticklabel style={/pgf/number format/fixed, /pgf/number format/precision=2},
  enlarge x limits=0.02,
  legend style={font=\small},
]
% Poisson(3)
\addplot[draw=sourceblue, fill=sourceblue!55] coordinates {
(0,0.0498)(1,0.1494)(2,0.2240)(3,0.2240)(4,0.1680)(5,0.1008)(6,0.0504)(7,0.0216)(8,0.0081)(9,0.0027)(10,0.0008)};
% Poisson(5)
\addplot[draw=quotientorange, fill=quotientorange!45] coordinates {
(0,0.0067)(1,0.0337)(2,0.0842)(3,0.1404)(4,0.1755)(5,0.1755)(6,0.1462)(7,0.1044)(8,0.0653)(9,0.0363)(10,0.0181)(11,0.0082)(12,0.0034)(13,0.0013)};
% Poisson(8) = sum
\addplot[draw=densityred, fill=densityred!40] coordinates {
(0,0.0003)(1,0.0027)(2,0.0107)(3,0.0286)(4,0.0573)(5,0.0916)(6,0.1221)(7,0.1396)(8,0.1396)(9,0.1241)(10,0.0993)(11,0.0722)(12,0.0481)(13,0.0296)(14,0.0169)(15,0.0090)(16,0.0045)(17,0.0021)(18,0.0009)};
\legend{$\Pois(3)$, $\Pois(5)$, sum $=\Pois(8)$}
\end{axis}
\end{tikzpicture}
\caption{The red bars are the distribution of $X+Y$ where $X\sim\Pois(3)$ and
$Y\sim\Pois(5)$ independently. It is exactly $\Pois(8)$. The PGF proof
(\ref{ch18:ex-pois-sum}) says this with no arithmetic: the exponents
$\lambda(s-1)$ and $\mu(s-1)$ add.}
\label{ch18:fig-pois-sum}
\end{figure}

\subsection{Reading moments off the PGF}

The PGF does double duty: it not only turns independent sums into products, it
also yields moments when you differentiate it. Differentiate the series term by term:
\[
  G_X'(s) = \sum_{k\ge 1} k\,p_k\, s^{k-1},
  \qquad
  G_X''(s) = \sum_{k\ge 2} k(k-1)\,p_k\, s^{k-2}.
\]
Now evaluate at $s=1$. The first gives $G_X'(1)=\sum_k k\,p_k = \E[X]$. The
second gives $G_X''(1)=\sum_k k(k-1)p_k = \E[X(X-1)] = \E[X^2]-\E[X]$. These
quantities $\E[X(X-1)\cdots(X-m+1)]$ are called \emph{factorial moments} (averages
of falling powers of $X$), and the PGF is their natural home. To recover the ordinary
variance, solve for the second moment $\E[X^2] = G_X''(1) + G_X'(1)$, then subtract
the squared mean $(\E[X])^2 = (G_X'(1))^2$.

\begin{proposition}[Mean and variance from the PGF]
\label{ch18:prop-pgf-moments}
If $X$ has PGF $G_X$, then
\[
  \E[X] = G_X'(1),
  \qquad
  \Var(X) = G_X''(1) + G_X'(1) - \big(G_X'(1)\big)^2 .
\]
\end{proposition}

\begin{example}[Poisson moments, instantly]
For $\Pois(\lambda)$ we have $G_X(s)=e^{\lambda(s-1)}$, so $G_X'(s)=\lambda
e^{\lambda(s-1)}$ and $G_X''(s)=\lambda^2 e^{\lambda(s-1)}$. At $s=1$ the
exponential is $1$, giving $G_X'(1)=\lambda$ and $G_X''(1)=\lambda^2$. Hence
$\E[X]=\lambda$ and
\[
  \Var(X) = \lambda^2 + \lambda - \lambda^2 = \lambda .
\]
The famous fact that a Poisson's mean and variance coincide drops out of one
exponential and two derivatives.
\end{example}

\begin{walkbox}
Use the PGF $G(s)=p/(1-qs)$ of the Geometric$(p)$ (failures before first
success) to show $\E[X]=q/p$.
\emph{Hint:} differentiate $p(1-qs)^{-1}$ to get $G'(s)=pq(1-qs)^{-2}$, then set
$s=1$ and remember $1-q=p$. You should find $G'(1)=pq/p^2 = q/p$. For practice,
push on to $G''(1)=2pq^2/p^3 = 2q^2/p^2$ and assemble
$\Var(X)=q/p^2$.
\end{walkbox}

\begin{applicationbox}[Branching processes and extinction]
PGFs are the native language of \emph{branching processes}, the model behind
nuclear chain reactions, surname survival, and epidemic spread. Suppose a single
ancestor starts generation $0$. In generation $1$, this individual produces a
random number of offspring with PGF $G(s)$. In generation $2$, each offspring
independently produces their own children using the same distribution $G(s)$.
A random sum of independent random variables corresponds to substituting one
PGF into another: the size of generation $2$ has PGF $G(G(s))$. By generation
$n$, the population size has PGF given by $n$-fold composition:
$G_n(s) = G(G(\cdots G(s)\cdots))$.

What is the probability of \emph{extinction} --- the event that the population
eventually drops to zero and stays there forever? If $d$ is the extinction
probability, conditioning on the number of offspring in the very first generation
shows that $d = G(d)$. The extinction probability is the smallest solution to
the fixed-point equation $G(s)=s$ in $[0,1]$. A whole theory --- including the
clean threshold that extinction is certain if and only if the mean offspring
count $\E[X]=G'(1)$ is $\le 1$ (excluding the trivial case $\P(X=1)=1$) --- falls
out of iterating one single-variable function. Composition of PGFs tracks
generational branching just as multiplication tracks independent sums.
\end{applicationbox}

\section{The Moment Generating Function}
\label{ch18:mgf}

The PGF has a built-in limitation: it is tailored to nonnegative integers,
where $s^X$ forms an ordinary power series. Real life is full of continuous
quantities --- waiting times, measurement errors, asset returns --- and for
those we need a transform that works whether $X$ is discrete, continuous,
positive, or negative. The fix is to replace the base $s$ by $e^t$. Then
$s^X = (e^t)^X = e^{tX}$ is well defined for \emph{any} real $X$.

\begin{definition}[Moment generating function]
\label{ch18:def-mgf}
The \emph{moment generating function} of a random variable $X$ is
\[
  M_X(t) \;=\; \E\!\left[e^{tX}\right],
\]
defined for those real $t$ where the expectation is finite. We say the MGF
\emph{exists} if $M_X(t)<\infty$ for all $t$ in some open interval
$(-\delta,\delta)$ around $0$.
\end{definition}

\begin{recoverybox}
\textbf{Mental check: the anchor of an MGF.}
What must $M_X(0)$ equal for \emph{any} random variable $X$?
Evaluate directly: $M_X(0) = \E[e^{0\cdot X}] = \E[1] = 1$.
Every legitimate MGF must pass through the point $(0,1)$, just as every PGF passes through $(1,1)$. If an MGF calculation gives $M_X(0) \ne 1$, an algebra error has happened along the way.
\end{recoverybox}

You can think of the MGF as the PGF with $s=e^t$ substituted in: where they both
exist, $M_X(t)=G_X(e^t)$. Everything good about the PGF carries over, and the
domain widens to all of $\R$.

\begin{theorem}[MGF of an independent sum]
\label{ch18:thm-mgf-product}
If $X$ and $Y$ are independent, then $M_{X+Y}(t)=M_X(t)\,M_Y(t)$, wherever both
sides are finite.
\end{theorem}

The proof is the same one line as before:
$M_{X+Y}(t)=\E[e^{t(X+Y)}]=\E[e^{tX}e^{tY}]=\E[e^{tX}]\E[e^{tY}]$, the middle
step by independence. Sums of independent variables become products of MGFs,
now for continuous variables too.

\subsection{Why it is called \emph{moment} generating}

Here is where the name is earned. Expand $e^{tX}$ in its Taylor series and take
the expectation term by term:
\[
  M_X(t) = \E\!\left[e^{tX}\right]
         = \E\!\left[\sum_{n=0}^\infty \frac{(tX)^n}{n!}\right]
         = \sum_{n=0}^\infty \frac{t^n}{n!}\,\E[X^n].
\]
Look at what this says. The MGF, expanded as a power series in $t$, has the
$n$th moment $\E[X^n]$ sitting on the coefficient of $t^n/n!$. The whole
sequence of moments is packed into one function, exactly the way the whole pmf
was packed into the PGF. To pull out the $n$th moment, differentiate $n$ times
and set $t=0$:

\begin{proposition}[Moments by differentiation]
\label{ch18:prop-mgf-moments}
If $M_X$ exists in a neighborhood of $0$, then for every $n\ge 1$,
\[
  \E[X^n] \;=\; M_X^{(n)}(0),
\]
the $n$th derivative of the MGF evaluated at $t=0$. In particular
$\E[X]=M_X'(0)$ and $\Var(X)=M_X''(0)-\big(M_X'(0)\big)^2$.
\end{proposition}

\begin{whybox}
Differentiating $\sum_n \E[X^n]\,t^n/n!$ once knocks the series down to
$\sum_{n\ge 1}\E[X^n]\,t^{n-1}/(n-1)!$, whose value at $t=0$ is the $n=1$ term,
$\E[X]$. Differentiating $n$ times and setting $t=0$ similarly isolates the
$\E[X^n]$ term --- every lower term has been differentiated to zero, every
higher term still carries a factor of $t$. The factorials in the series are
placed exactly so the $n$th derivative cleanly recovers $\E[X^n]$ with no stray
constants. The MGF is a moment dispenser, and $t=0$ is the lever.
\end{whybox}

\begin{example}[The exponential distribution]
\label{ch18:ex-expo-mgf}
Let $X\sim\Expo(\lambda)$ with density $f(x)=\lambda e^{-\lambda x}$ for
$x\ge 0$. Then
\[
  M_X(t) = \int_0^\infty e^{tx}\,\lambda e^{-\lambda x}\,\dd x
         = \lambda\int_0^\infty e^{-(\lambda-t)x}\,\dd x
         = \frac{\lambda}{\lambda - t},
  \qquad t<\lambda .
\]
The integral converges only when $\lambda - t>0$, which is why the MGF lives on
the half-line $t<\lambda$ and not all of $\R$ --- a first hint that MGFs can be
fragile. To read off moments, expand as a geometric series:
\[
  M_X(t) = \frac{1}{1 - t/\lambda} = \sum_{n=0}^\infty \frac{t^n}{\lambda^n}
         = \sum_{n=0}^\infty \frac{n!}{\lambda^n}\cdot\frac{t^n}{n!}.
\]
Matching against $\sum_n \E[X^n]\,t^n/n!$ gives $\E[X^n]=n!/\lambda^n$ for all
$n$ at once. In particular $\E[X]=1/\lambda$ and $\E[X^2]=2/\lambda^2$, so
$\Var(X)=2/\lambda^2 - 1/\lambda^2 = 1/\lambda^2$. One series, all the moments.
\end{example}

\begin{example}[The normal distribution]
\label{ch18:ex-normal-mgf}
The MGF of the standard normal $Z\sim\Normal(0,1)$ is the keystone of the CLT
proof, so it is worth doing carefully. With density
$\varphi(z)=\frac{1}{\sqrt{2\pi}}e^{-z^2/2}$,
\[
  M_Z(t) = \frac{1}{\sqrt{2\pi}}\int_{-\infty}^\infty e^{tz}\,e^{-z^2/2}\,\dd z .
\]
Complete the square in the exponent: $tz - \tfrac12 z^2 = -\tfrac12(z-t)^2 +
\tfrac12 t^2$. Pull the $t$-only piece out of the integral:
\[
  M_Z(t) = e^{t^2/2}\cdot
           \frac{1}{\sqrt{2\pi}}\int_{-\infty}^\infty e^{-(z-t)^2/2}\,\dd z
         = e^{t^2/2}\cdot 1 = e^{t^2/2},
\]
since the remaining integral is just a normal density (centered at $t$)
integrating to $1$. So $M_Z(t)=e^{t^2/2}$, finite for \emph{all} real $t$ ---
the normal's tails are so thin that $e^{tZ}$ has finite expectation no matter
how large $t$ is. For a general $X\sim\Normal(\mu,\sigma^2)$, write
$X=\mu+\sigma Z$ and use the scaling rule below to get
\[
  M_X(t) = e^{\mu t + \frac12 \sigma^2 t^2}.
\]
\end{example}

\subsection{How the MGF responds to shifting and scaling}

Two small rules make the MGF a pleasure to manipulate.

\begin{proposition}[Affine transformations]
\label{ch18:prop-affine}
For constants $a,b$, $\;M_{aX+b}(t) = e^{bt}\,M_X(at).$
\end{proposition}

\begin{whybox}
$M_{aX+b}(t)=\E[e^{t(aX+b)}]=\E[e^{(at)X}\,e^{bt}] = e^{bt}\,\E[e^{(at)X}] =
e^{bt}M_X(at)$. The constant $b$ slides outside as the factor $e^{bt}$; the
multiplier $a$ reaches inside and rescales the argument. This is exactly how you
turn the standard normal's MGF into the general normal's: with $a=\sigma$,
$b=\mu$, $\;M_{\mu+\sigma Z}(t)=e^{\mu t}M_Z(\sigma t)=e^{\mu t}e^{\sigma^2
t^2/2}$, recovering Example~\ref{ch18:ex-normal-mgf}.
\end{whybox}

\begin{example}[Sum of independent normals, by transform]
\label{ch18:ex-normal-sum}
Let $X\sim\Normal(\mu_1,\sigma_1^2)$ and $Y\sim\Normal(\mu_2,\sigma_2^2)$ be
independent. Multiply their MGFs:
\[
  M_{X+Y}(t)=e^{\mu_1 t+\frac12\sigma_1^2 t^2}\,
             e^{\mu_2 t+\frac12\sigma_2^2 t^2}
            =e^{(\mu_1+\mu_2)t+\frac12(\sigma_1^2+\sigma_2^2)t^2}.
\]
That is the MGF of a $\Normal(\mu_1+\mu_2,\ \sigma_1^2+\sigma_2^2)$. So the sum
of independent normals is normal, with means and variances adding --- a fact
that costs a serious convolution integral to prove directly, but here is just
``add the exponents.'' The normal family is \emph{closed under addition}, and
the MGF makes that visible at a glance.
\end{example}

\begin{trapbox}
The product rule needs \textbf{independence}. If $X$ and $Y$ are dependent,
$M_{X+Y}(t)$ is generally \emph{not} $M_X(t)M_Y(t)$ --- the step
$\E[e^{tX}e^{tY}]=\E[e^{tX}]\E[e^{tY}]$ fails. A common slip is to ``add
variances'' or ``multiply MGFs'' for correlated variables. Both are independence
facts. For dependent sums you must keep the joint distribution (or use the
\emph{joint} MGF $\E[e^{t_1 X + t_2 Y}]$ and its mixed partials).
\end{trapbox}

\begin{applicationbox}[Chernoff bounds and tail behavior]
The MGF is not only for identifying distributions; it is the engine behind the
sharpest elementary tail bounds in all of probability. For any $t>0$, Markov's
inequality (the fundamental bound $\P(Y\ge c)\le \E[Y]/c$ for a nonnegative
random variable $Y$, which we explore deeply in Chapter~19) applied to $e^{tX}$ gives
\[
  \P(X\ge a) = \P\!\left(e^{tX}\ge e^{ta}\right)
             \le \frac{\E[e^{tX}]}{e^{ta}}
             = e^{-ta}\,M_X(t).
\]
Because this holds for \emph{every} $t>0$, you may optimize over $t$: the
\emph{Chernoff bound} is $\P(X\ge a)\le \min_{t>0} e^{-ta}M_X(t)$. For an
independent sum $S_n=X_1+\cdots+X_n$, the MGF factorizes:
$M_{S_n}(t) = \prod_{i=1}^n M_{X_i}(t)$. For i.i.d.\ variables this becomes
$(M_{X_1}(t))^n$, so the bound on the tail of the average $\overline{X}_n = S_n/n$
takes the form $\big[e^{-ta} M_{X_1}(t)\big]^n$ --- a quantity that shrinks
\emph{exponentially fast} as the sample size $n$ grows. This exponential decay
is the mathematical backbone of generalization bounds in statistical learning,
the performance of randomized hashing algorithms, and concentration inequalities.
Whenever a machine-learning paper writes ``with probability at least $1-\delta$,''
an optimized MGF Chernoff bound is almost certainly doing the heavy lifting underneath.
\end{applicationbox}

\section{Uniqueness: Why the Transform Is Trustworthy}
\label{ch18:uniqueness}

Every trick in this chapter ends with the phrase ``\dots and therefore $X+Y$ is
Poisson'' or ``\dots and therefore $X+Y$ is normal.'' Each such conclusion
relies on one essential fact: that the MGF \emph{determines} the
distribution. If two random variables have the same MGF, they have the same
distribution --- so recognizing the MGF lets you name the distribution. Without
this, computing $M_{X+Y}$ would be a dead end; with it, recognition is proof.

\begin{theorem}[Uniqueness of the MGF]
\label{ch18:thm-uniqueness}
Suppose $X$ and $Y$ have MGFs that are finite and equal on some open interval
$(-\delta,\delta)$ around $0$: $M_X(t)=M_Y(t)$ for all $\abs{t}<\delta$. Then
$X$ and $Y$ have the same distribution.
\end{theorem}

We will not prove this in full --- the real proof routes through the
characteristic function and Fourier inversion, which we meet next. But the
\emph{reason} it is plausible is worth stating, because it explains the
underlying mechanism.

\begin{whybox}
A finite MGF on an interval around $0$ is real-analytic there, and its Taylor
coefficients are the moments $\E[X^n]/n!$. So equal MGFs near $0$ means
\emph{all moments agree}. The remaining content of the theorem is that, when the
MGF exists in a neighborhood of $0$, the moments pin the distribution down
uniquely. Intuitively, knowing every moment is knowing the average of every
power $X,X^2,X^3,\dots$, hence the average of every polynomial, hence
(by approximating indicators with polynomials) the probability of every event.
The existence of the MGF is what guarantees the tails are thin enough for this
program to close --- which is also exactly the fine print, addressed in the
trap below.
\end{whybox}

\begin{trapbox}[Moments alone are \emph{not} always enough]
\label{ch18:trap-moments}
It is tempting to believe ``same moments $\Rightarrow$ same distribution''
unconditionally. It is false. There exist genuinely \emph{different}
distributions with \emph{identical} moments of every order. The classic example
is the lognormal distribution: if $\log X\sim\Normal(0,1)$, one can perturb its
density by a wiggle that leaves every moment $\E[X^n]$ unchanged, producing a
different distribution with exactly the same moment sequence. The catch is that
the lognormal's moments grow so fast (its MGF is infinite for every $t>0$) that
they fail to determine the distribution. The escape hatch is precisely the
hypothesis of Theorem~\ref{ch18:thm-uniqueness}: if the \emph{MGF} is finite in
an open interval around $0$, the tails are controlled, the moment series
converges, and uniqueness holds. This is why we are careful to say ``the MGF
exists,'' not merely ``the moments exist.''
\end{trapbox}

\begin{connectionbox}[The Stieltjes moment problem: when moments fail to pin down a distribution]
The question of whether a sequence of moments $\E[X^n]$ uniquely determines a
distribution occupied Pafnuty Chebyshev, Andrei Markov, and Thomas Stieltjes in the
late nineteenth century. Stieltjes (1894) gave the first systematic treatment,
inventing the Stieltjes integral (integrating directly against a general distribution
function that may possess both discrete jumps and smooth continuous curves) specifically
to solve it. Later, Hans Hamburger and Felix Hausdorff pinned down exact conditions on
the moment sequence. The practical rule of thumb is Carleman's condition: a distribution
on the real line is uniquely determined by its moments if the series
$\sum_{k=1}^\infty (\E[X^{2k}])^{-1/(2k)}$ diverges (meaning high-order moments do
not grow faster than $(2k)!$). Finiteness of the MGF in an open interval around $0$
guarantees Carleman's condition with plenty of room to spare. The lognormal
distribution's failure to be pinned down by its moments was a landmark historical
warning: even an infinite list of summary numbers can leave a distribution ambiguous
if the tails are heavy.
\end{connectionbox}

\begin{example}[Recognition in practice]
Suppose you compute, for some $X$ built out of a complicated experiment, that
$M_X(t)=(1-2t)^{-3}$ for $t<1/2$. You do not need to know how $X$ was built. You
recall that a $\Gam(\alpha,\beta)$ with rate $\beta$ has MGF
$(1-t/\beta)^{-\alpha}$. Matching, $\beta=1/2$ and $\alpha=3$. By uniqueness, $X
\sim \Gam(3, 1/2)$ --- equivalently a chi-squared with $6$ degrees of freedom.
The MGF acted as a fingerprint: identify the fingerprint, identify the
distribution.
\end{example}

\section{Characteristic Functions: The Transform That Always Exists}
\label{ch18:cf}

The MGF has one serious limitation:
\emph{it may not exist.} We saw the exponential's MGF diverge for $t\ge\lambda$, and
the lognormal's diverge for every $t>0$. Heavy-tailed distributions --- the Cauchy
distribution being the headline case --- have $\E[e^{tX}]=\infty$ for every
$t\ne 0$, so their MGF is useless. Since the MGF is built from $e^{tX}$, which
explodes when $X$ is large and $t>0$, any distribution with fat enough tails
breaks it.

The repair is one of the most elegant moves in analysis: put an imaginary unit
in the exponent. Replace $e^{tX}$ by $e^{itX}$, where $i=\sqrt{-1}$.

\begin{definition}[Characteristic function]
\label{ch18:def-cf}
The \emph{characteristic function} of a random variable $X$ is
\[
  \varphi_X(t) \;=\; \E\!\left[e^{itX}\right]
              \;=\; \E[\cos(tX)] + i\,\E[\sin(tX)],
  \qquad t\in\R .
\]
\end{definition}

\begin{happeningbox}
Why does the $i$ rescue everything? Because $e^{itX}$ does not grow. By Euler's
formula, $e^{itX}=\cos(tX)+i\sin(tX)$, a point on the unit circle in the complex
plane. Its magnitude is $\abs{e^{itX}}=1$ for every value of $X$, no matter how
huge. Where the real exponential $e^{tX}$ shoots off to infinity for large $X$,
the complex exponential just \emph{spins} around the circle (Figure~\ref{ch18:fig-complex-rotation}). So
$\abs{\varphi_X(t)}=\abs{\E[e^{itX}]}\le \E\abs{e^{itX}} = \E[1] = 1 < \infty$
\emph{always}. The integral defining $\varphi_X$ converges for every
distribution on Earth, because we are averaging bounded quantities. The CF is
the MGF with its dangerous growth tamed into harmless rotation.
\end{happeningbox}

\begin{figure}[H]
\centering
\begin{tikzpicture}[scale=0.95]
  % Panel 1: Real MGF growth explosion
  \begin{scope}[shift={(0,0)}]
    \draw[draw=gray!30, fill=gray!3, rounded corners=4pt] (-1.0,-1.55) rectangle (5.0,3.9);
    \node[anchor=north, font=\small\bfseries, densityred!90!black] at (2.0,3.75) {1. Real MGF: Explosive Growth $e^{tx}$};

    % Axes
    \draw[->, thick, gray!70] (-0.6,0) -- (4.4,0) node[right, font=\scriptsize] {$x$};
    \draw[->, thick, gray!70] (0,-0.2) -- (0,2.7) node[above, font=\scriptsize] {$y$};

    % Exploding exponential curve e^{0.8 x}
    \draw[densityred, very thick]
      (-0.5, 0.67) -- (0, 1.0) -- (0.5, 1.49) -- (1.0, 2.22) -- (1.2, 2.55)
      node[above right=-2pt, font=\scriptsize\bfseries] {$y = e^{tx}$ ($t > 0$)};

    % Divergence callout for heavy tails
    \node[align=center, font=\scriptsize, draw=densityred!70, fill=white, rounded corners=2pt, inner sep=2pt]
      at (2.0, -0.75) {For heavy tails (e.g.\ Cauchy):\\[1pt]
      $f(x) \sim \frac{1}{x^2} \implies \displaystyle\int_{-\infty}^\infty e^{tx} f(x)\,\dd x = \infty$\\[1pt]\scriptsize (MGF fails to exist)};
  \end{scope}

  % Arrow connecting panels
  \draw[->, >=stealth, very thick, structurepurple] (5.2, 1.0) -- (6.3, 1.0)
    node[midway, above, font=\scriptsize\bfseries, align=center] {Multiply\\by $i$};

  % Panel 2: Complex CF bounded unit circle rotation
  \begin{scope}[shift={(8.8,0)}]
    \draw[draw=gray!30, fill=gray!3, rounded corners=4pt] (-2.2,-1.55) rectangle (3.8,3.9);
    \node[anchor=north, font=\small\bfseries, sourceblue!90!black] at (0.8,3.75) {2. Complex CF: Bounded Unit Circle};

    % Axes: Re and Im
    \draw[->, thick, gray!70] (-1.8,0.8) -- (3.2,0.8) node[right, font=\scriptsize] {$\mathrm{Re}$};
    \draw[->, thick, gray!70] (0.8,-0.8) -- (0.8,2.6) node[above, font=\scriptsize] {$\mathrm{Im}$};

    % Center of circle is at (0.8, 0.8), radius = 1.4
    \coordinate (C) at (0.8, 0.8);
    \draw[structurepurple!70, thick] (C) circle (1.4);
    \node[structurepurple!80!black, font=\tiny\bfseries] at ($(C)+(1.2, -1.1)$) {$|z| = 1$};

    % Points spinning on the unit circle
    \foreach \ang/\lbl in {35/{x_1}, 110/{x_2}, 205/{x_3}, 290/{x_4}} {
      \coordinate (P\lbl) at ($(C)+(\ang:1.4)$);
      \fill[distblue] (P\lbl) circle (1.6pt);
      \draw[gray!50, thin] (C) -- (P\lbl);
    }
    \node[font=\tiny\bfseries, distblue] at ($(C)+(35:1.68)$) {$e^{itx_1}$};
    \node[font=\tiny\bfseries, distblue] at ($(C)+(110:1.68)$) {$e^{itx_2}$};

    % Center of mass vector average phi_X(t)
    \coordinate (Phi) at ($(C)+(15:0.65)$);
    \draw[->, very thick, sourceblue!90!black] (C) -- (Phi);
    \node[font=\scriptsize\bfseries, sourceblue!90!black] at (1.55, 0.45) {$\varphi_X(t) = \E[e^{itX}]$};
    \fill[sourceblue!90!black] (Phi) circle (2pt);

    \node[align=center, font=\scriptsize, draw=sourceblue!70, fill=white, rounded corners=2pt, inner sep=2pt]
      at (0.8, -1.08) {Trapped in unit disk: $|\varphi_X(t)| \le 1$\\[1pt]\scriptsize Always finite for every distribution};
  \end{scope}
\end{tikzpicture}
\caption{Why the characteristic function always exists. (Left) The real exponential $e^{tx}$ diverges toward $+\infty$ on the real line, causing the MGF integral $\int e^{tx} f(x)\,\dd x$ to blow up for heavy-tailed distributions. (Right) The complex exponential $e^{itx} = \cos(tx) + i\sin(tx)$ rotates around the unit circle without growing ($|e^{itx}|=1$). The characteristic function $\varphi_X(t) = \E[e^{itX}]$ is the center-of-mass vector average of these points, guaranteed to remain trapped inside the unit disk $|\varphi_X(t)| \le 1$.}
\label{ch18:fig-complex-rotation}
\end{figure}

\begin{theorem}[Properties of the characteristic function]
\label{ch18:thm-cf-props}
For every random variable $X$:
\begin{enumerate}
  \item \emph{(Always exists, bounded.)} $\varphi_X(t)$ is defined for all
        $t\in\R$ and $\abs{\varphi_X(t)}\le 1$, with $\varphi_X(0)=1$.
  \item \emph{(Sums.)} If $X\indep Y$, then
        $\varphi_{X+Y}(t)=\varphi_X(t)\,\varphi_Y(t)$.
  \item \emph{(Affine.)} $\varphi_{aX+b}(t)=e^{ibt}\varphi_X(at)$.
  \item \emph{(Moments, when they exist.)} If $\E\abs{X}^n<\infty$, then
        $\varphi_X$ is $n$ times differentiable and
        $\varphi_X^{(n)}(0)=i^n\,\E[X^n]$.
  \item \emph{(Uniqueness.)} If $\varphi_X=\varphi_Y$ as functions on $\R$, then
        $X$ and $Y$ have the same distribution.
\end{enumerate}
\end{theorem}

Properties (2) and (3) have the identical one-line proofs as for the MGF ---
nothing changed but a letter $i$.

\begin{whybox}[Moments from the characteristic function]
Why does $\varphi_X^{(n)}(0) = i^n \E[X^n]$? Expand the complex exponential in its Taylor series:
\[
  e^{itX} \;=\; \sum_{k=0}^n \frac{(itX)^k}{k!} \;+\; R_n(tX)
          \;=\; 1 + itX - \frac{t^2}{2}X^2 + \cdots + \frac{i^n t^n}{n!}X^n + R_n(tX).
\]
Taking expectations term by term (rigorously justified whenever $\E\abs{X}^n < \infty$):
\[
  \varphi_X(t) \;=\; 1 + it\,\E[X] - \frac{t^2}{2}\,\E[X^2] + \cdots + \frac{i^n t^n}{n!}\,\E[X^n] + o(t^n).
\]
Differentiating $n$ times with respect to $t$ and evaluating at $t=0$ isolates the $n$th moment:
\[
  \varphi_X'(0) = i\,\E[X], \qquad \varphi_X''(0) = i^2\,\E[X^2] = -\E[X^2], \qquad \dots, \qquad \varphi_X^{(n)}(0) = i^n\,\E[X^n].
\]
Hence the mean is $\E[X] = \frac{1}{i}\varphi_X'(0) = -i\varphi_X'(0)$, and the variance is
$\Var(X) = \E[X^2] - (\E[X])^2 = -\varphi_X''(0) + \big(\varphi_X'(0)\big)^2$.
\end{whybox}

Property (5), uniqueness, is the deep one, and unlike
the MGF version it comes with \emph{no fine print}: because the CF always
exists, it \emph{always} determines the distribution. There is even an explicit
recipe to get the distribution back.

\begin{theorem}[Fourier inversion]
\label{ch18:thm-inversion}
If $X$ has a density $f$ and $\varphi_X$ is integrable, then
\[
  f(x) \;=\; \frac{1}{2\pi}\int_{-\infty}^\infty e^{-itx}\,\varphi_X(t)\,\dd t .
\]
\end{theorem}

This formula captures the main idea of the chapter: the
characteristic function is the \emph{Fourier transform of the density}, and the
density is the \emph{inverse Fourier transform of the characteristic function}.
Nothing is lost in the round trip. The transform is not a lossy summary like the
mean or the variance; it is an exact re-encoding of the entire distribution,
and the inversion integral is the proof.

\begin{connectionbox}[This is Fourier analysis]
If you have met Fourier series or transforms (in the companion calculus book, or
in any signal-processing course), the characteristic function is familiar:
writing $\varphi_X(t)=\int e^{itx}f(x)\,\dd x$, it is precisely the
Fourier transform of the density $f$ (with a sign and $2\pi$ convention chosen
by probabilists). The convolution theorem of Fourier analysis --- ``the
transform of a convolution is the product of the transforms'' --- is exactly our
Theorem~\ref{ch18:thm-cf-props}(2): the density of $X+Y$ is the convolution
$f_X * f_Y$, and its transform is the product $\varphi_X\varphi_Y$. The
machinery of this chapter applies this core tool of
harmonic analysis to distributions. The reason transforms turn convolution into multiplication
is the same reason here as there: $e^{it(x+y)}=e^{itx}e^{ity}$.
\end{connectionbox}

\begin{paradoxbox}[Averaging Cauchys gets you nowhere]
\label{ch18:ex-cauchy}
The standard Cauchy density is $f(x)=\frac{1}{\pi(1+x^2)}$. Its tails are so
heavy that $\E\abs{X}=\infty$: the mean does not even exist, let alone the MGF.
Yet its characteristic function is perfectly well behaved:
\[
  \varphi_X(t) = e^{-\abs{t}} .
\]
(The derivation is a contour integral; take the result on faith for now.) Watch
what the product rule does. If $X_1,\dots,X_n$ are independent standard Cauchys,
then $\overline{X}_n=\frac1n\sum X_i$ has CF
\[
  \varphi_{\overline X_n}(t)
   = \prod_{k=1}^n \varphi_{X_k}(t/n)
   = \big(e^{-\abs{t}/n}\big)^n
   = e^{-\abs{t}} = \varphi_X(t).
\]
By uniqueness, $\overline{X}_n$ has the \emph{same} distribution as a single
$X$. Averaging a thousand independent Cauchy observations gives you something
distributed exactly like one observation --- the sample mean does not
concentrate, and the Law of Large Numbers fails. This famous failure is
invisible to the MGF (which does not exist here) but laid bare by the CF in
three lines. It also demonstrates why the CF is the tool of
choice: it works where everything else has given up. The resolution of the
``paradox'' is that the LLN requires a finite mean, which the Cauchy
lacks; the heavy tails mean any average is dominated by its largest term.
\end{paradoxbox}

\begin{figure}[H]
\centering
\begin{tikzpicture}
\begin{axis}[
  width=12.5cm, height=5.6cm,
  domain=-4:4, samples=160,
  xlabel={$t$}, ylabel={$\varphi_X(t)$},
  title={Characteristic functions are bounded by $1$ and peak at $t=0$},
  title style={font=\small},
  ymin=-0.25, ymax=1.08,
  legend style={font=\small, at={(0.98,0.97)}, anchor=north east},
  axis lines=left,
]
\addplot[sourceblue, very thick] {exp(-x*x/2)};
\addplot[densityred, very thick] {exp(-abs(x))};
\addplot[quotientorange, very thick, dashed] {sin(deg(x))/x};
\legend{{$\Normal(0,1)$: $e^{-t^2/2}$}, {Cauchy: $e^{-|t|}$}, {$\Unif(-1,1)$: $\tfrac{\sin t}{t}$}}
\end{axis}
\end{tikzpicture}
\caption{Three characteristic functions. All start at $\varphi(0)=1$ and stay
within $\pm 1$ --- a guarantee that holds for every distribution. The normal's
CF is again a Gaussian bell; the Cauchy's is a sharp tent $e^{-\abs t}$; the
uniform's is the oscillating $\sin(t)/t$, dipping below zero, which a real MGF
could never do.}
\label{ch18:fig-cf}
\end{figure}

\begin{aestheticbox}
The three transforms form a natural hierarchy. The PGF handles counts and turns
convolution into polynomial multiplication. The MGF widens the door to all real
random variables and dispenses moments by differentiation --- but it can crash
on heavy tails. The CF puts an $i$ in the exponent, trading explosive growth for
harmless rotation on the unit circle, and thereby \emph{always} exists, always
determines the distribution, and comes with an exact inversion formula. Three
transforms, one idea, increasing robustness. Each is the right tool for a
different job, and together they convert hard questions about adding randomness
into easy questions about multiplying functions.
\end{aestheticbox}

\section{A Preview of the Central Limit Theorem}
\label{ch18:clt-preview}

We close with the result that makes this entire chapter feel inevitable. The
Central Limit Theorem says that the sum (suitably standardized) of many
independent, identically distributed pieces is approximately normal, almost
regardless of what the pieces look like. The proof, which we give in full in the
CLT chapter, is essentially a one-page computation \emph{with characteristic
functions} --- and now you can see its skeleton.

\begin{focusbox}
\textbf{The CLT in transform language.} Let $X_1,X_2,\dots$ be i.i.d.\ with mean
$0$ and variance $1$ (subtract the mean and divide by the SD to arrange this),
and let $\varphi$ be their common characteristic function. Form the standardized
sum
\[
  S_n = \frac{X_1+\cdots+X_n}{\sqrt n}.
\]
By the sum rule (\ref{ch18:thm-cf-props}, part 2) and the affine rule (part 3),
\[
  \varphi_{S_n}(t)
   = \prod_{k=1}^n \varphi\!\left(\frac{t}{\sqrt n}\right)
   = \left[\varphi\!\left(\frac{t}{\sqrt n}\right)\right]^n .
\]
Now Taylor-expand $\varphi$ near $0$ using the moment property: since $\E[X]=0$
and $\E[X^2]=1$,
\[
  \varphi(u) = 1 + i\,\E[X]\,u - \tfrac12 \E[X^2]\,u^2 + o(u^2)
             = 1 - \tfrac12 u^2 + o(u^2).
\]
Substitute $u=t/\sqrt n$:
\[
  \varphi_{S_n}(t)
   = \left[1 - \frac{t^2}{2n} + o\!\left(\frac{1}{n}\right)\right]^n
   \xrightarrow[n\to\infty]{} e^{-t^2/2},
\]
using the limit $(1+a_n/n)^n\to e^a$ when $a_n\to a$. But $e^{-t^2/2}$ is
exactly the characteristic function of the standard normal
(Example~\ref{ch18:ex-normal-mgf}, with the same calculation under the $i$). By
uniqueness and a continuity theorem, $S_n$ converges in distribution to
$\Normal(0,1)$. \emph{That is the Central Limit Theorem.}
\end{focusbox}

% Underlying X: symmetric Bernoulli +-1, CF cos(t). S_n CF = [cos(t/sqrt n)]^n.
\begin{figure}[H]
\centering
\begin{tikzpicture}
\begin{axis}[
  width=12.5cm, height=5.6cm,
  domain=-3:3, samples=160,
  xlabel={$t$}, ylabel={$\varphi_{S_n}(t)$},
  title={The CF of a standardized sum climbs toward the Gaussian $e^{-t^2/2}$},
  title style={font=\small},
  ymin=-0.2, ymax=1.08,
  xtick={-3,-2,-1,0,1,2,3},
  ticklabel style={font=\small},
  legend style={font=\small, at={(0.98,0.97)}, anchor=north east},
  axis lines=left,
]
\addplot[quotientorange, very thick] {cos(deg(x))};
\addplot[structurepurple, very thick, dashed] {(cos(deg(x/1.7320508)))^3};
\addplot[targetgreen, very thick, dotted] {(cos(deg(x/3.1622777)))^10};
\addplot[densityred, very thick] {exp(-x*x/2)};
\legend{{$n=1$: $\cos t$}, {$n=3$}, {$n=10$}, {limit $e^{-t^2/2}$}}
\end{axis}
\end{tikzpicture}
\caption{Take the $X_i$ to be $\pm1$ fair coin flips (CF $\cos t$). As $n$ grows,
the CF of the standardized sum $S_n=(X_1+\cdots+X_n)/\sqrt n$, namely
$[\cos(t/\sqrt n)]^n$, climbs toward the standard-normal CF $e^{-t^2/2}$
(solid red). This is the Central Limit Theorem watched in transform space:
multiplication of CFs, then a limit.}
\label{ch18:fig-clt}
\end{figure}

Pause on how much heavy lifting the transform just did. The distribution of a
sum of $n$ convolved-together pieces would be a nightmare to track directly. In
transform space it is a single power $[\varphi(t/\sqrt n)]^n$ (watch it converge
in Figure~\ref{ch18:fig-clt}), the limit is a one-line Taylor expansion, and the
answer $e^{-t^2/2}$ recognizes itself as the normal. \emph{Why is the bell curve universal?} Because near zero every
characteristic function looks like $1-\frac12(\text{variance})u^2$, and that
quadratic --- the only surviving term after centering sets the mean to zero and
scaling fixes the variance --- exponentiates into the Gaussian. The normal is universal because
$-\frac12 t^2$ is what is left when you keep only the first two moments. The
generating function is what let us see it.

\begin{historybox}
The story of generating functions and transforms is the story of probability
transforming from an ad-hoc collection of gambling calculations into a rigorous
branch of modern analysis. Abraham de Moivre in the 1730s first used generating
polynomials to solve the gambler's ruin problem. In 1812, Pierre-Simon Laplace
greatly expanded the idea in his \emph{Th\'eorie analytique des probabilit\'es},
using formal integral transforms to discover the bell-curve limit for sums of
errors. Laplace and his contemporaries were stuck in a computational swamp:
evaluating $n$-fold convolution integrals by hand quickly became impossible.
Transforming functions gave them an operational calculus to turn convolutions
into algebra, but they lacked modern analysis and treated complex exponentials
and divergent series as formal tricks. Nineteenth-century mathematicians viewed
these manipulations with deep suspicion as lacking rigorous foundations. Controversy deepened
when Thomas Stieltjes in 1894 showed that two distinct probability distributions
could share identical moments of every order (such as the lognormal), revealing
that the method of moments could fail to determine a distribution. The breakthrough
came in the 1920s when Paul L\'evy and Aleksandr Khinchin built the rigorous
theory of characteristic functions on the complex unit circle. L\'evy proved
the celebrated \emph{continuity theorem} (if characteristic functions converge
pointwise to a continuous function, the underlying distributions converge),
providing the first watertight proof of the Central Limit Theorem. Today,
while one-dimensional transforms are fully solved by Fourier inversion,
inverting transforms in hundreds or thousands of dimensions remains computationally
formidable due to the curse of dimensionality (the exponential explosion of grid
points). Researchers in machine-learning diffusion models, quantum tomography,
and medical imaging still actively work on high-dimensional transform inversion
and phase retrieval (reconstructing a probability distribution when only incomplete
transform data can be observed).
\end{historybox}

\section{Summary}
\label{ch18:summary}

Three transforms, one principle. Each encodes a whole distribution in a single
function so that adding independent random variables becomes multiplying their
transforms, and moments become derivatives.

\begin{center}
\begin{tabular}{@{}lccc@{}}
\toprule
 & \textbf{PGF} & \textbf{MGF} & \textbf{CF} \\
\midrule
Definition       & $\E[s^X]$ & $\E[e^{tX}]$ & $\E[e^{itX}]$ \\
For which $X$    & $\N$-valued & needs thin tails & \emph{any} $X$ \\
Always exists?   & on $\abs s\le1$ & not always & \textbf{yes} \\
Independent sum  & multiply & multiply & multiply \\
Moments          & $G^{(n)}(1)$ (factorial) & $M^{(n)}(0)=\E[X^n]$
                 & $\varphi^{(n)}(0)=i^n\E[X^n]$ \\
Determines dist.? & yes & yes (if exists) & \textbf{always} \\
\bottomrule
\end{tabular}
\end{center}

\begin{itemize}
  \item The \textbf{PGF} $G_X(s)=\E[s^X]$ hangs the pmf on powers of $s$.
        Convolution becomes polynomial multiplication; factorial moments come
        from derivatives at $s=1$.
  \item The \textbf{MGF} $M_X(t)=\E[e^{tX}]=G_X(e^t)$ extends the idea to all
        real $X$. Its Taylor coefficients are the moments
        ($\E[X^n]=M^{(n)}(0)$); it determines the distribution \emph{when it
        exists in an interval around $0$}.
  \item The \textbf{CF} $\varphi_X(t)=\E[e^{itX}]$ puts an $i$ in the exponent,
        bounding $\abs{e^{itX}}=1$ so the transform always exists. It is the
        Fourier transform of the density, always determines the distribution,
        and is the engine of the CLT proof.
\end{itemize}

The single fact under all of it: $\;e^{(\cdot)(X+Y)}=e^{(\cdot)X}e^{(\cdot)Y}$,
so independence turns ``sum of variables'' into ``product of transforms.'' That
one identity is why generating functions exist, why convolutions simplify, and
why the normal distribution is universal.

\subsection*{Looking ahead: inequalities and concentration}
Generating functions provide exact distributions and moments when closed forms exist.
In many applications, however, only partial information is available---such as the mean,
variance, or bounded range of a variable.
Chapter~19 turns from exact calculations to \textbf{inequalities and concentration},
showing how Markov's inequality, Chebyshev's inequality, and Chernoff bounds
(using the MGF) provide exponential tail bounds on deviations from the mean.

\section{Exercises}
\label{ch18:exercises}

\begin{exercise}[Warm-up PGFs]
\label{ch18:ex-warmup}
Write the probability generating function of each, and check that it equals $1$
at $s=1$.
\begin{enumerate}[label=(\alph*)]
  \item A single fair four-sided die, values $\set{1,2,3,4}$.
  \item $X\sim\Bin(3,1/2)$, by direct summation, and confirm it matches
        $(\tfrac12+\tfrac12 s)^3$.
  \item The constant random variable $X=7$.
\end{enumerate}
\end{exercise}

\begin{exercise}[Sum of dice via PGF]
\label{ch18:ex-twodice}
Let $G(s)=\frac16(s+s^2+s^3+s^4+s^5+s^6)$ be the PGF of one fair die. The sum of
two independent dice has PGF $G(s)^2$. Expand $G(s)^2$ and read off
$\P(\text{sum}=7)$ as the coefficient of $s^7$. Confirm it equals $6/36$.
\end{exercise}

\begin{exercise}[Negative binomial closure]
\label{ch18:ex-negbin}
The Geometric$(p)$ (failures before first success) has PGF $p/(1-qs)$. Using the
product rule, find the PGF of a sum of $r$ independent Geometric$(p)$ variables.
Identify the resulting distribution (it is Negative Binomial). What is its mean,
via $G'(1)$?
\end{exercise}

\begin{exercise}[MGF moment extraction]
\label{ch18:ex-mgf-moments}
Let $X\sim\Expo(\lambda)$ with $M_X(t)=\lambda/(\lambda-t)$ for $t<\lambda$.
\begin{enumerate}[label=(\alph*)]
  \item Compute $M_X'(t)$ and $M_X''(t)$ and evaluate at $t=0$ to recover
        $\E[X]=1/\lambda$ and $\E[X^2]=2/\lambda^2$.
  \item Use these to confirm $\Var(X)=1/\lambda^2$.
  \item From the series $\sum_n (n!/\lambda^n)\,t^n/n!$, write down $\E[X^5]$
        with no further calculus.
\end{enumerate}
\end{exercise}

\begin{exercise}[Identify by MGF]
\label{ch18:ex-identify}
A random variable $Y$ is found to have $M_Y(t)=e^{3(e^t-1)}$ for all $t$. By
recognizing this MGF (substitute $s=e^t$ into a PGF you know), name the
distribution of $Y$ and state its mean and variance. Then explain in one
sentence which theorem of this chapter licenses the conclusion.
\end{exercise}

\begin{exercise}[Sum of independent normals]
\label{ch18:ex-normal-sum-prob}
Let $X\sim\Normal(1,4)$ and $Y\sim\Normal(-2,9)$ be independent. Using
$M_{\Normal(\mu,\sigma^2)}(t)=e^{\mu t+\frac12\sigma^2 t^2}$ and the product
rule, find the MGF of $2X - Y + 5$, and hence its distribution. (Use the affine
rule on each piece first.)
\end{exercise}

\begin{exercise}[Why the $i$ matters]
\label{ch18:ex-cauchy-mgf}
For the standard Cauchy density $f(x)=1/[\pi(1+x^2)]$:
\begin{enumerate}[label=(\alph*)]
  \item Show that $\E[e^{tX}]=\infty$ for every $t\ne 0$, so the MGF does not
        exist. (Compare the integrand's tail $e^{tx}/x^2$ to a divergent
        integral for $t>0$; handle $t<0$ by symmetry.)
  \item Given $\varphi_X(t)=e^{-\abs t}$, use the product and affine rules to
        find the characteristic function of $\overline X_n=\frac1n\sum_{i=1}^n
        X_i$ for i.i.d.\ standard Cauchys, and conclude that $\overline X_n$ has
        the same distribution as a single $X_i$.
\end{enumerate}
\end{exercise}

\begin{exercise}[CF basics]
\label{ch18:ex-cf-basics}
\begin{enumerate}[label=(\alph*)]
  \item Show directly from the definition that $\abs{\varphi_X(t)}\le 1$ and
        $\varphi_X(0)=1$ for every random variable.
  \item Show that $\varphi_X(-t)=\overline{\varphi_X(t)}$ (complex conjugate),
        and deduce that if $X$ is symmetric about $0$ (i.e.\ $X$ and $-X$ have
        the same distribution) then $\varphi_X(t)$ is real-valued.
\end{enumerate}
\end{exercise}

\begin{exercise}[Uniqueness, carefully]
\label{ch18:ex-uniqueness}
Two students argue. Student A says: ``If $X$ and $Y$ have all the same moments,
they have the same distribution.'' Student B says: ``Not necessarily.'' Who is
right, and under exactly what extra hypothesis does Student A's claim become
true? Name the standard counterexample and explain in one or two sentences why
its moments fail to determine it.
\end{exercise}

\begin{exercise}[CLT skeleton, filled in]
\label{ch18:ex-clt}
Let $X_1,\dots,X_n$ be i.i.d.\ with mean $0$, variance $1$, and characteristic
function $\varphi$. Starting from $\varphi(u)=1-\tfrac12 u^2 + o(u^2)$, fill in
every step showing that the CF of $S_n=(X_1+\cdots+X_n)/\sqrt n$ converges to
$e^{-t^2/2}$ as $n\to\infty$. State precisely where you use (i) independence,
(ii) the affine rule, (iii) the second-moment Taylor expansion, and (iv) the
limit $(1+a_n/n)^n\to e^a$.
\end{exercise}

\begin{exercise}[Branching extinction]
\label{ch18:ex-branching}
In a branching process each individual independently has $0$, $1$, or $2$
offspring with probabilities $\tfrac14,\tfrac14,\tfrac12$.
\begin{enumerate}[label=(\alph*)]
  \item Write the offspring PGF $G(s)$ and compute the mean offspring count
        $G'(1)$. Is it above or below $1$?
  \item The extinction probability is the smallest root in $[0,1]$ of
        $G(s)=s$. Solve the resulting quadratic and report the extinction
        probability. Does it match your expectation from part (a)?
\end{enumerate}
\end{exercise}

\begin{exercise}[Looking ahead: Chernoff bound for coin flips]
\label{ch18:ex-chernoff-coin}
Let $X_1,\dots,X_n$ be independent Bernoulli$(1/2)$ variables (fair coin flips
taking values in $\set{0,1}$), and let $S_n = \sum_{i=1}^n X_i$ be the total
number of heads.
\begin{enumerate}[label=(\alph*)]
  \item Show that the MGF of a single flip is $M_{X_1}(t) = \frac{1+e^t}{2}$,
        and deduce that $M_{S_n}(t) = \big(\frac{1+e^t}{2}\big)^n$.
  \item Use the elementary calculus inequality $1+u \le e^u$ (with $u=(e^t-1)/2$)
        to show that $\frac{1+e^t}{2} \le \exp\big(\frac{e^t-1}{2}\big)$.
  \item Using Markov's inequality, for any threshold $a > n/2$ and any $t > 0$,
        show that
        \[
          \P(S_n \ge a) \;\le\; e^{-ta}\left(\frac{1+e^t}{2}\right)^n .
        \]
  \item Notice that this bound has the form $(c(t))^n$ for a constant $c(t) < 1$
        when $t$ is chosen smartly. Explain in one sentence why this guarantees
        that the probability of getting mostly heads decays \emph{exponentially}
        with the number of flips, previewing the concentration inequalities of
        Chapter~19.
\end{enumerate}
\end{exercise}
